%% ma: L10-B13-0073
%% lop: 10
%% chuong: 5
%% bai: 13
%% dang: 10.13.15
%% loai: TL
%% muc: TH
%% dap_an: "-"
%% nguon: "chuyen-de-cac-so-dac-trung-cua-mau-so-lieu-khong-ghep-nhom-toan-10-chuong-trinh-moi.pdf (GV. Phan Nhật Linh) (P03), Bài 02 Dạng 2 Bài tập tự luận 10"
%% kien_thuc: "Lớp 10 Bài 13 (số trung bình, trung vị, tứ phân vị, mốt)"
%% trang_thai: chinh-thuc
%% ghi_chu: "Tự duyệt theo yêu cầu giáo viên 22/10/2026."
\begin{ex}
Số lượng học sinh giỏi quốc gia năm học 2018 -- 2019 của $10$ trường THPT được cho như sau: $0$; $0$; $4$; $0$; $0$; $0$; $10$; $0$; $6$; $0$.

a) Tìm số trung bình, mốt và các tứ phân vị của mẫu số liệu.

b) Giải thích tại sao tứ phân vị thứ nhất và trung vị trùng nhau.
\loigiai{a) Sắp xếp: $0$; $0$; $0$; $0$; $0$; $0$; $0$; $4$; $6$; $10$. $\overline x=\dfrac{20}{10}=2$; mốt $0$; $Q_2=\dfrac{0+0}2=0$; $Q_1=0$; $Q_3=4$.

b) Có $7$ trong $10$ giá trị bằng $0$, nên cả giá trị thứ $3$ (vị trí của $Q_1$) và hai giá trị thứ $5$, thứ $6$ (vị trí của $Q_2$) đều bằng $0$.}
\end{ex}
